\subsection{Vector Space of Linear Maps}

Very light on the notes side, probably because this is such a critical new topic. Problems 11 and 17 were basically impossible, I attempted them but it was not possible for me to solve them unfortunately.


\begin{enumerate}
    \item Suppose $b, c \in \mathbb{R}$. Define $T \colon \mathbb{R}^3 \to \mathbb{R}^2$ by
$T(x, y, z) = (2x - 4y + 3z + b, 6x + cxyz)$.
Show that $T$ is linear if and only if $b = c = 0$.

    \begin{proof}
        Suppose that $b \neq 0$. Then, $T(0) = (b, 0) \neq 0$, so its not a linear map. 

        Suppose that $c \neq 0$. Then, look at $T((1, 1, 1) + (1, 1, 1)) = T((2, 2, 2)) = (2, 12 + 8c) \neq T(1, 1, 1) + T(1, 1, 1) = (2, 12 + c)$.

        Now let's show that when $b = 0, c = 0$, $T$ is a linear map. Suppose we have $a, b, \in \R^3$. 

        \begin{align*}
            T(a + b) &= (2a_1 + 2b_1 - 4a_2 - 4b_2 + 3a_3 + 3b_3, 6a_1 + 6b_1) \\
            &= (2a_1 - 4a_2 + 3a_3, 6a_1) + (2b_1 - 4b_2 + 3b_3, 6b_1) \\
            &= T(a) + T(b) 
        \end{align*}
    \end{proof}


    \item Suppose $b, c \in \mathbb{R}$. Define $T \colon \mathcal{P}(\mathbb{R}) \to \mathbb{R}^2$ by
    
$$Tp = \left(3p(4) + 5p'(6) + bp(1)p(2), \int_{-1}^{2} x^3 p(x) \, dx + c \sin p(0)\right)$$

Show that $T$ is linear if and only if $b = c = 0$.

\begin{proof}
    This thing is so bad honestly its easier to prove that the sum of a linear map and a non-linear map is non-linear. 

    Suppose we have functions $f \in \mathcal L(U, W)$ and $g: U \to W$ (in this case, $U$ and $W$ are just sets who happen to be interpretable as vector spaces). If $f + g$ was a linear map, so would $f + g - f$, which clearly isn't.

    Then, we know that linear maps are preserved under addition, scalar multiplication, and function composition, so everything except the terms with $b$ and $c$ are cleared. $bp(1)p(2)$ is not a linear map (when $b \neq 0$) because it does not preserve scalar multiplication. Multiplying $p$ by a constant multiplites $bp(1)p(2)$ by its square. $c\sin p(0)$ clearly violates both addition and multiplication ($\sin$ isn't linear!). So, $b = c = 0$. Since everything else is a linear map, $T$ must be too.
\end{proof}

\item Suppose that $T\in\mathcal L(\mathbb F^n,\mathbb F^m)$. Show that there exist scalars $A_{j,k}\in\mathbb F$ for
\hspace{1em} $j=1,\dots,m$ and $k=1,\dots,n$ such that for every
\hspace{1em} $(x_1,\dots,x_n)\in\mathbb F^n$,
\[
T(x_1,\dots,x_n)=\bigl(\nobreak\sum_{k=1}^n A_{1,k}x_k,\dots,\nobreak\sum_{k=1}^n A_{m,k}x_k\bigr).
\]

\begin{proof}
    Consider the standard basis for $\F^n$. The Linear Map Lemma states that a linear map is entirely determined by where it sends its basis vectors. Let 
    
    $$Te_i = (A_{1, i}, \dots, A_{m, i})$$
    
    Then, 

    \begin{align*}
        T(x_1, \dots, x_n) &= Tx_1 + \dots + Tx_n \\
        &= (A_{1,1}, \dots, A_{m, 1}) + \dots + (A_{1, n}, \dots A_{m, n}) \\
        &= \bigl(\nobreak\sum_{k=1}^n A_{1,k}x_k,\dots,\nobreak\sum_{k=1}^n A_{m,k}x_k\bigr)
    \end{align*}
\end{proof}

\item[5] Prove that $\mathcal L(V, W)$ is a vector space, as was asserted in 3.6.

\begin{proof}
    We are given the following definitions for addition and scalar multiplication:

    Given $S, T, \in \mathcal L(V, W)$ and $ \lambda \in \F$,

    \begin{align*}
        (S + T)(v) = Sv + Tv \\
        (\lambda T)(v) = \lambda(Tv)
    \end{align*}

    Just like the proofs for vector spaces themselves, we can strap our properties to the existing properties of our ring. 

    \begin{enumerate}
        \item Communativity
        
        \begin{align*}
            (S + T)(v) &= Sv + Tv \\
            &= Tv + Sv \text{ (Communativity of $W$)} \\
            &= (T + S)(v)
        \end{align*}

        \item Associativity 
        
        \begin{align*}
            ((S + T) + H)(v) &= (S + T)(v) + H(v) \\
            &= Sv + Tv + Hv \\
            &= Sv + (Tv + Hv) \\
            &= (S + (T + H))(v)
        \end{align*}

        \item Additive Inverse
        
        Define $(-T)(v) = -Tv$.

        \begin{align*}
            (T + (-T))(v) = Tv - Tv = 0
        \end{align*}

        \item Multiplicative Identity
        
        Our identity element is 1. This is trivially the identity, shush.

        \item Distributive property
        
        \begin{align*}
            (\lambda(S + T))(v) &= \lambda(S + T)(v) \\
            &= \lambda(Sv + Tv) \\
            &= \lambda Sv + \lambda Tv
            &= (\lambda S + \lambda T)(v)
        \end{align*}

    \end{enumerate}
\end{proof}

\item[8] Give an example of a function $\varphi\colon \mathbb{R}^2 \to \mathbb{R}$ such that
\[
\varphi(av) = a\varphi(v)
\]
for all $a \in \mathbb{R}$ and all $v \in \mathbb{R}^2$, but $\varphi$ is not linear.

\begin{proof}
    Maybe something like 

    $$f(a, b) = \begin{cases}
        a & a=b \\
        0 & a \neq b
    \end{cases}$$ 

    would work.
\end{proof}

\item[9] Give an example of a function $\varphi\colon \mathbb{C} \to \mathbb{C}$ such that
\[
\varphi(w+z)=\varphi(w)+\varphi(z)
\]
for all $w,z\in\mathbb{C}$ but $\varphi$ is not linear. (Here $\mathbb{C}$ is thought
of as a complex vector space.)

\begin{proof}
    Something like $\varphi(z) = \bar{z}$ would work. It satisfies additivity but 

    $$i\overline{1 + i} \neq \overline{i(1 + i)}$$
\end{proof}

\item[10] Prove or give a counterexample: If $q \in \mathcal{P}(\mathbb{R})$ and $T \colon \mathcal{P}(\mathbb{R}) \to \mathcal{P}(\mathbb{R})$ is
defined by $Tp = q \circ p$, then $T$ is a linear map.

\begin{proof}
    Not true in general, since something like $x^2$ violates the additivity constraint. ($(1 + 1)^2 \neq 2^2$)
\end{proof}


\item[11] Suppose $V$ is finite-dimensional and $T \in \mathcal L(V)$. Prove that $T$ is a scalar
multiple of the identity if and only if $ST = TS$ for every $S \in \mathcal L(V)$.

\begin{proof}
    I cannot seem to find a good solution to this exercise online, so it is skipped. Maybe I will come back to it.
\end{proof}

\item[12] Suppose $U$ is a subspace of $V$ with $U \neq V$. Suppose $S \in \mathcal{L}(U, W)$ and
$S \neq 0$ (which means that $Su \neq 0$ for some $u \in U$). Define $T \colon V \to W$ by
\[
Tv=
\begin{cases}
Sv & \text{if } v \in U,\\
0 & \text{if } v \in V \text{ and } v \notin U.
\end{cases}
\]
Prove that $T$ is not a linear map on $V$.

\begin{proof}
    Suppose there was some $u \in U$ where $Su \neq 0$. Then, we have some $v \in V, v \notin U$. $u + v \notin U$, so $T(u + v) = 0$. But also $T(u + v) = Tu + Tv = Tu \neq 0$. Contradiction!
\end{proof}

\item[15] Suppose $v_1,\dots,v_m$ is a linearly dependent list of vectors in $V$. Suppose
also that $W\neq\{0\}$. Prove that there exist $w_1,\dots,w_m\in W$ such that no
$T\in\mathcal L(V,W)$ satisfies $Tv_k=w_k$ for each $k=1,\dots,m$.

\begin{proof}
    We have some linear combination, $a_i$ not all 0, such that 

    $$a_1v_1 + \dots + a_mv_m = 0$$

    Let $a_i$ be the first non-zero index. We let $w_i$ be an arbitrary non-zero vector in $W$ and all other $j \neq i , w_j = 0$. 

    Then, we have that 

    $$T(a_1v_1 + \dots + a_mv_m) = 0$$ 

    But also 
    \begin{align*}
        T(a_1v_1 + \dots + a_mv_m) &= a_1T(v_1) + \dots + a_mT(v_m) \\
        &= a_1w_1 + \dots + a_mw_m \\
        &= a_iw_i \\
        &\neq 0
    \end{align*}

    So no such $T$ exists.
    
\end{proof}

\item[17] Suppose $V$ is finite-dimensional. Show that the only two-sided ideals of
$\mathcal{L}(V)$ are $\{0\}$ and $\mathcal{L}(V)$.
A subspace $\mathcal{E}$ of $\mathcal{L}(V)$ is called a two-sided ideal of $\mathcal{L}(V)$ if $TE \in \mathcal{E}$ and
$ET \in \mathcal{E}$ for all $E \in \mathcal{E}$ and all $T \in \mathcal{L}(V)$.

\begin{proof}
    This problem is cooked.
\end{proof}

\end{enumerate}